\section{Introduction}\label{sec:introduction}
Vision-language-action (VLA) models connect visual observations and language instructions to robot commands. Their pretrained representations offer a starting point for manipulation, but their action interface depends on the embodiment used during training. When a new robot requires different action dimensions or control conventions, adapting the model can involve replacing or reinitializing its action head or expert. In this setting, useful visual and language representations may be available before the policy can reliably execute even a basic movement.

Learning a complete instruction directly from binary success is difficult when the initial policy rarely succeeds. Consider ``put the apple on the plate'': the robot must first approach the apple, grasp it, lift it, carry it, and release it correctly. A failed rollout gives the same final outcome whether the robot never reaches the object or drops it just before placement. Training each component separately can make learning easier, but those components must subsequently work together from a normal task start under one instruction.

Our approach follows this structure. We first use dense rewards to learn individual skills, giving the policy feedback before it can complete them. We then transition to binary skill-success rewards and self-imitation of successful experience. Finally, we train the complete task with staged rewards so that earlier skills and later task progress contribute to learning. The aim is to adapt VLA-based control to a new embodiment through a practical training sequence, rather than to make a particular robot architecture the subject of the method.

We investigate embodiment adaptation using SmolVLA and a cable-driven parallel robot (CDPR), whose gripper is positioned by coordinated cable actuation. This platform serves as a demanding example of a robot with a different mechanical and control structure from conventional serial arms. Its role is to test adaptation under a substantial embodiment change; the approach does not rely on a cable-specific task objective. The evaluated implementation learns a new residual action module while retaining the pretrained SmolVLA action expert and representations (Sec.~\ref{sec:method}).

A complementary experiment applies phase-aware feedback to an already trained OpenVLA-OFT manipulator on LIBERO. It examines whether intermediate task structure is also useful after an action policy is available. This study uses its own robot, scenes, and training setup, with no CDPR evaluation. It tests one component of the approach, rather than repeating the complete new-embodiment training sequence.

The paper contributes a dense-to-binary skill-learning schedule with self-imitation and staged task composition; an implementation for adapting VLA-based control to a cable-driven parallel manipulator; and evaluations of complete-task execution and stage-aware feedback. The intended scope is manipulation across robot embodiments. Evidence from these experiments establishes the tested cases, while broader transfer remains an empirical question.
